\documentclass[11pt]{article}
\usepackage[margin=1in]{geometry}
\usepackage{amsmath,amssymb,amsthm}
\usepackage[T1]{fontenc}
\usepackage{lmodern}
\usepackage{microtype}
\usepackage[hidelinks]{hyperref}
\newtheorem{theorem}{Theorem}
\newcommand{\T}{\mathbb{T}}
\newcommand{\N}{\mathbb{N}}
\newcommand{\C}{\mathbb{C}}
\newcommand{\R}{\mathbb{R}}
\newcommand{\norm}[1]{\left\lVert#1\right\rVert}
\title{Proof of Conjecture 00000000934\\[4pt]
\large An unconditionally convergent series in complex $L^1(\R/\mathbb{Z})$}
\author{C0ldSmi1e}
\date{October 9, 2026}
\begin{document}
\maketitle

\begin{abstract}
We construct a series in complex-valued $L^1$ of the circle that converges
unconditionally after every termwise sign choice and admits no rearrangement
that converges absolutely in the $L^1$ norm. The construction uses the Fourier
monomials with harmonic coefficients. In fact, arbitrary complex unit phases
may replace signs. The accompanying Lean 4 project proves the complete
existence assertion with the usual almost-everywhere quotient and norm.
\end{abstract}

\section{Statement and conventions}
The English source asserts that there is an unconditionally convergent series
in $L^1(\T)$ which remains unconditionally convergent after every sign change
and admits no absolute rearrangement. The exact English and Chinese source
is preserved in \texttt{ORIGINAL.md}.

We use the complex scalar convention: $\T=\R/\mathbb{Z}$ has normalized Haar
measure $\mu$, with $\mu(\T)=1$. For this circumference-one circle, $\mu$ is
also the usual Lebesgue volume. The space $L^1=L^1(\T;\C)$ consists of
integrable functions modulo equality almost everywhere, with
$\norm{f}_1=\int_{\T}|f|\,d\mu$.
Unconditional convergence means convergence in this norm of the net of sums
over all finite subsets of the index set. A sign choice is any function
$\varepsilon:\N\to\{-1,1\}$, with no support or eventual-constancy restriction.
An absolute rearrangement would be a bijection $\pi:\N\to\N$ for which
$\sum_n\norm{f_{\pi(n)}}_1<\infty$. Thus ``absolute'' refers to the norm of
the stated Banach space, rather than an alternative pointwise convergence condition.

The source's parenthetical reference to the boundary of Orlicz's theorem is
retained as descriptive attribution. It specifies no additional endpoint
parameter or sharpness assertion. We make no separate historical or endpoint claim.

\begin{theorem}
There exists $f:\N\to L^1(\T;\C)$ such that $\sum_n f_n$ is
unconditionally convergent, every infinite termwise scalar sign choice also
gives an unconditionally convergent series, and every permutation fails to
give an absolutely norm-convergent series.
\end{theorem}

\section{Construction and unconditional convergence}
For $n\geq1$, let
\[
 e_n(x)=\exp(2\pi i n x),\qquad f_n=\frac{e_n}{n},
 \qquad f_0=0.
\]
The exponential is invariant under integer translations of $x$, so it is
well-defined on $\T$. Each term is bounded and integrable. The initial zero
term only permits indexing by $\N$; it does not affect any convergence assertion.

The functions $e_n$ are orthonormal in $L^2(\T;\C)$:
\[
 \int_{\T}\overline{e_m}e_n\,d\mu=
 \begin{cases}1,&m=n,\\0,&m\ne n.\end{cases}
\]
For distinct integers $m,n$, the integrand is the character $e_{n-m}$.
Translation by $1/(2(n-m))$ multiplies it by $-1$, while Haar invariance
preserves its integral; the integral is therefore zero. For $m=n$ the
integrand is identically one.

Fix an arbitrary sequence $\varepsilon_n\in\C$ with $|\varepsilon_n|=1$.
For finite sets $F,G\subset\N$, orthogonality gives
\[
 \left\lVert\sum_{n\in F}\varepsilon_n f_n
       -\sum_{n\in G}\varepsilon_n f_n\right\rVert_2^2
   =\sum_{\substack{n\in F\mathbin{\triangle}G\\n\geq1}}\frac1{n^2}.
\]
Given $\delta>0$, convergence of $\sum_{n\geq1}n^{-2}$ gives $M$ for which
the tail above $M$ is less than $\delta^2$. If both $F$ and $G$ contain
$\{0,\ldots,M\}$, the displayed difference has norm less than $\delta$.
The finite-subset net is Cauchy. Completeness of $L^2$ yields its limit.

For every $g\in L^2$, Cauchy--Schwarz and $\mu(\T)=1$ imply
\[
 \norm{g}_1=\int_{\T}|g|\cdot1\,d\mu
       \leq\left(\int_{\T}|g|^2\,d\mu\right)^{1/2}
       =\norm{g}_2.
\]
The identity on almost-everywhere classes is consequently a bounded linear
map $L^2\to L^1$. It sends the convergent finite-subset net to a convergent
net in the actual $L^1$ norm. Hence $\sum_n\varepsilon_n f_n$ converges
unconditionally for every unit-phase sequence. The original series follows
by taking all phases equal to $1$; every requested sign choice is included.
The limiting element may depend on the phase sequence.

For any fixed phase sequence, every bijection $\pi$ gives the same norm
limit: its finite prefix sets $\{\pi(0),\ldots,\pi(N-1)\}$ eventually contain
each prescribed finite subset of $\N$. This also explains directly the
relation between the finite-subset and permutation formulations.

\section{Failure of every absolute rearrangement}
The pointwise equality $|e_n(x)|=1$ and $\mu(\T)=1$ give
\[
 \norm{f_n}_1=\frac1n\quad(n\geq1),\qquad\norm{f_0}_1=0.
\]
Unit phases preserve these norms. The harmonic series diverges, since each
block $2^k\leq n<2^{k+1}$ contributes at least $1/2$.
A bijection of indices does not change the supremum of finite sums of a
nonnegative family. Therefore for every unit-phase sequence and every
permutation $\pi$,
\[
 \sum_{n=0}^{\infty}\norm{\varepsilon_{\pi(n)}f_{\pi(n)}}_1=+\infty.
\]
In particular no absolute rearrangement of the original series exists.
Together with unconditional convergence above, this proves all clauses of
the theorem for a single sequence.\hfill$\square$

\section{Formal correspondence and verification}
The project pins Lean 4.19.0 and Mathlib v4.19.0. All declarations below are
in the namespace \texttt{Conjecture934} of \texttt{Solution.lean}.

\begin{center}
\begin{tabular}{p{0.42\textwidth}p{0.49\textwidth}}
\hline
Mathematical content & Lean declarations\\\hline
Actual circle, measure, a.e. quotient and norm &
\texttt{Circle}, \texttt{\textmu}, \texttt{L1}, \texttt{L2},
\texttt{measure\_eq\_volume}, \texttt{L1\_norm\_integral}\\[4pt]
Explicit Fourier sequence and representatives &
\texttt{term}, \texttt{term\_ae}, \texttt{representative\_on\_real}\\[4pt]
Bounded identity inclusion $L^2\to L^1$ &
\texttt{inclusionLinear\_exists}, \texttt{inclusionLinear\_apply},
\texttt{inclusion\_norm\_le}, \texttt{inclusion\_exists},
\texttt{inclusion\_apply}, \texttt{inclusion\_fourier}\\[4pt]
Every infinite phase and sign choice &
\texttt{phase\_L2\_summable}, \texttt{phase\_summable},
\texttt{IsSignChoice}, \texttt{sign\_summable}\\[4pt]
Explicit norm limits and all permutations &
\texttt{phase\_unconditional\_norm},
\texttt{phase\_permutation\_norm}\\[4pt]
Harmonic norm and no absolute ordering &
\texttt{term\_norm}, \texttt{no\_absolute\_rearrangement},
\texttt{phase\_norm\_partial\_sums\_diverge}\\[4pt]
Complete existential source assertion & \texttt{conjecture}\\[4pt]
Stronger assertion for every unit phase & \texttt{strengthened\_conjecture}\\\hline
\end{tabular}
\end{center}

Mathlib's \texttt{Summable} means existence of a limit for the net of finite
subset sums. Here it uses the norm topology of \texttt{MeasureTheory.Lp}.
The proof invokes the library's orthonormal Fourier family and the
orthogonal-family square-summability theorem, then transports convergence
through the explicitly established inclusion. Chosen inclusion maps are
selected from proved existence statements; pointwise action lemmas fully
fix their values on every input. They introduce no additional assumptions.

The frozen mathematical core passes a fresh source replay with warnings
treated as errors. Its audit inspects all 36 owned declarations, including
generated helpers, and 30,946 transitive declarations. Only the standard
axioms \texttt{propext}, \texttt{Classical.choice}, and \texttt{Quot.sound}
occur; the audit accepts no unsafe declaration or owned axiom.
No numerical computation is needed for the mathematics. The auxiliary
programs check provenance and proof dependencies.

\texttt{PROOF.md} gives the full declaration mapping;
\texttt{VERIFICATION.md} retains both successful checks and resolved
development failures. The final package's \texttt{VERIFICATION-PACKAGE.md}
provides the reviewer workflow and additional verification records.
Local compilation and independent review do not constitute maintainer
acceptance. This submission claims the complex-valued theorem stated above,
without a separate formalized real-valued or Orlicz-endpoint assertion.
\end{document}
